% Source PDF page 10
\begin{frame}[t,allowframebreaks]{Structure}
  \fontsize{15}{18}\selectfont
  \begin{itemize}
    \item Many \textcolor{CommandGreen}{\bfseries iptables} commands have the following structure:
  \end{itemize}
  \vspace{0.08cm}
  {\color{CommandGreen}\ttfamily\fontsize{13.5}{16.5}\selectfont \#iptables [-t <table-name>]\par}
  <command><chain-name> <parameter-1> <option-1> <parameter-n><option-n>\par
  \begin{itemize}
    \item <table-name> : allows the user to select a table other than
  \end{itemize}
  the default \textcolor{CommandGreen}{\bfseries filter} table to use with the command.\par
  \begin{itemize}
    \item <command> : dictates a specific action to perform, such
  \end{itemize}
  as appending or deleting the rule specified by the <chain- name> option.\par
  \begin{itemize}
    \item <chain-name> :are pairs of parameters and options that
  \end{itemize}
  define what will happen when a packet matches the rule\par
  \vspace{0.08cm}
  {\color{CommandGreen}\ttfamily\fontsize{13.5}{16.5}\selectfont iptables -A INPUT -s 0/0 -i eth0 -d 192.168.1.1 -p TCP -j ACCEPT\par}
\end{frame}
